\section{Additional Scope Checks}
\label{sec:analysis}

\subsection{Descriptive variation across MMLU groups}

Across the standard broad MMLU groups, keep14 remains below the intact base and
ShortGPT remains above keep14. The realized keep14 scores range from .295 on
STEM to .379 on Other; ShortGPT improves STEM/humanities/social science/Other
by 11.0/13.6/21.1/19.7 points but does not close the base gap. The advantage is
therefore broad rather than driven by one group. The complete group and
57-subject inventories remain in Appendix~\ref{app:mmlu-subjects}.

\subsection{Construction effects beyond MMLU}

The complete profile in Table~\ref{tab:downstream} shows that ShortGPT's
advantage is not confined to MMLU. Relative to keep14 at displayed step 200k,
it gains 4.1 points on HellaSwag, 3.8 on ARC-Challenge, and 9.1 on BoolQ; its
core6 macro rises from .5938 to .6215. The gain is therefore broad, but highly
target-dependent: the 15.6-point letter-MMLU difference is much larger than the
2.8-point core6 difference, the 1.8-point content-MMLU difference, and the
1.7/3.6/0.7-point closed-book gains on PopQA/TriviaQA/NQ-open.

This profile strengthens the design lesson from the same-depth comparison.
Construction choices do not simply translate every benchmark by a constant
amount; they reshape the evaluation profile. Reporting one aggregate endpoint
would hide both ShortGPT's broad benefit and the unusually large
answer-letter-MMLU effect. Table~\ref{tab:downstream} reports all eleven
multiple-choice/likelihood tasks and all endpoint rows; the
closed-book columns are in Table~\ref{tab:main-results}.

\subsection{The shallow observations are not a depth law}
The retained shallow rows stop at 121k (keep8), 83.5k (keep10), and 124k (keep12), whereas keep14 reaches 200k. Their stopping points and FLOPs differ, so the ladder is an endpoint inventory rather than a matched-step or compute-normalized ordering.

Within keep8, several likelihood-style tasks and PPL improve over the observed
interval while aggregate MMLU remains near chance. Together with the keep14
trajectory, this identifies matched stopping-rule and compute-normalized depth
experiments as the next step.

\subsection{Scope checks do not constitute replications}
A more compressed OLMo-2-1B prefix+fresh-tail run shows falling PPL with
near-chance MMLU, and an existing Qwen3-8B endpoint is also near chance after
200k. The 1B arm changes scale and retained fraction; the Qwen arm additionally
changes architecture, corpus, available evaluations, and --- unlike every other
arm here --- starts from an instruction-tuned rather than a base checkpoint. We
include them as directional context only, not as matched replications or evidence
of cross-family universality (Appendix Table~\ref{tab:app-crossfamily}).

\subsection{Synthesis}
Across these scope checks, the same audit logic holds: a target score must be
interpreted together with its interface, construction, and training budget.
The main OLMo-2-7B study supports this conclusion directly; the shallow, 1B,
and Qwen observations indicate where matched replications would be most useful.
